\section{The lifting exponent and its denominator consequence}

% BEGIN A85D1
\paragraph{A85D1. A lift with exactly zero monodromy.}
Use the three realized modules of R85P1--3 and the exact triple
$M_0\xrightarrow F M_1\xrightarrow G M_2$. Denote the induced
operator on principal parts by $T=[E\nabla]$. Put
\[
 w=2(N-1)R,\qquad c_1=b(R+1),\qquad c_2=w+c_1,\qquad h_L=(N+1)b .
\]
Suppose $y'_2\in Y'_2$, $m_2\in M_2^\Delta$ and
$\mathrm{Mon}(y'_2+\varphi m_2)=0$. By E85C4 lift
$p^w y'_2$ and $p^w m_2$ separately to $y'_1$ and $m_1$ in
the corresponding integral submodules. Write
$x=y'_1+\varphi m_1$ and $\omega=\mathrm{Mon}(x)$. We have
$Gx=p^w(y'_2+\varphi m_2)$. By functoriality $G\omega=0$;
by M85N4--5, $p^{c_1}\omega\in Z_1$. Apply K85F2 with
$s=w$ to this class. There is $z_0\in Z_0$ satisfying
$Fz_0=p^{c_2}\omega$.

Since $x\in Y_1$, K85F4 puts $\omega$ in the image of $T$ on
$Z_{1,L}$. It also belongs to the subspace $V=F Z_{0,L}$, as
we see by dividing the last equality by $p^{c_2}$. This subspace
is $T$-stable by horizontality. M85N2 gives
$V\cap T Z_{1,L}=T V$, so some $q\in Z_{0,L}$ has
$FTq=\omega$. Injectivity in K85F2 gives $p^{c_2}Tq=z_0$.
All jumps of the kernel are in $U$ by K85F3 or R85P3.
M85N6 gives $y_0\in Y'_0$ with
$\mathrm{Mon}(y_0)=p^n Tq$ for $n=c_2+b+h_L$. The two vectors
\[
            p^n y'_1-Fy_0,\qquad p^n m_1
\]
belong to $Y'_1$ and $M_1^\Delta$. The monodromy of their
Frobenius sum is $p^n\omega-Fp^n Tq=0$. Its image under $G$
is exactly $p^{n+w}(y'_2+\varphi m_2)$ since $GF=0$. Thus the
middle input to Q84T4 holds with $k_*=c_1+2w+b+h_L$.
% END A85D1

% BEGIN A85D2
\paragraph{A85D2. The absolute admissible lifting constant on the original tests.}
Fix a test surjection $g:T_1\to T_2$ and an integral finite
cocycle $c$ of $T_2$ as in Q84T1. First extend the coefficients
and pass to the common tame field as in Q84T1--2, retaining the
full $G_{\mathbb Q_p}$ descent. R85P1--3 and E85C1 supply the
models for this test and for the extension attached to $c$.
P85C2 gives $y_2\in Y_2$ representing $p^w c$ with zero
monodromy. Write $p^b y_2=y'_2+m_2$ using
$m_2=E^b y_2$ and $p^b-E^b\in u^e A$. Then
$y'_2+\varphi m_2=p^b y_2+(\varphi-1)m_2$. This has the
same class as $p^b y_2$ modulo integral coboundaries and still
has zero monodromy. Apply A85D1. P85C3--5 return the resulting
zero-monodromy vector, multiplied by $p^{c_1}$, to a finite
integral cocycle of $T_1$ on the full original Galois group.
P85C6 compares its projection with $p^{d}c$ up to an integral
coboundary of $T_2$. The costs are
\[
 \begin{aligned}
 d&=w+b+k_*+c_1\\
  &=6(N-1)R+(N+5+2R)b\le 7\{NR+b(N+R)\}.
 \end{aligned}
\]
The last inequality uses $N\ge2$, $R\ge1$. Q84T2 supplies the
other conclusion of the same test, $4R\le N^2$, for the original
ratio period. The coefficient
retraction in Q84T1 descends the identity to the initial lattices.
All constants depend only on the displayed numerical invariants,
even for tests of unbounded rank. Raising the exponent to
$40\{NR+b(N+R)\}$ proves the native $\mathrm{admT}(40)$
under $40(HN^2+N^3)<p$ for the usual period data. The same
argument applies with any larger integral constant under its own
prime guard.
% END A85D2

% BEGIN A85D3
\paragraph{A85D3. Passing the type-prime bound to the uncut family.}
Take the family in P84B3 with its original published inputs, twist
condition, original Galois preconditions and four compatibility
conditions on the tested members. Fix the type-prime constant
$C_T=40$ before choosing a member or a place. At each degree-one
completion T84L1 identifies the local test with a surjection of the
kind used in A85D2. K84P1--4 and T84L5 supply the common tame
characters and the labelled jumps on its direct or dual support.
Thus A85D2 supplies the quantitative premise of P84B3 on every
such support, including all ranks tested there. T84L2--6 bound
the direct and dual local defects. P84B3 then gives the point
estimate on the chosen cut and the family assertion
$\mathrm{LinkN}$ after return to the original uncut family.
The finiteness, summability and scale tests are the original
ones in P84B2, with the cut chosen before the lift. Consequently
this family implication requires no further type-prime lifting
hypothesis. To obtain the selected obligation
$\mathrm{GLOBAL\_DENOMINATOR}$ one must still produce the
member hypotheses within its original family and design quantifiers.
% END A85D3
